% CS455-167
\section{Overview and final scope}

The Campus Ride-Pooling and Split-Fare System is a web application for IIT Kanpur students travelling
between campus and transport hubs in the city. A student offers a ride or requests a seat on someone
else's ride, the fare is shared equally among the people in the car, and the group coordinates through
a private chat. The application calculates who owes what but does not handle money: settlement happens
outside it, by cash or UPI, and is only recorded in the app.

The system runs at \url{https://campus-ridepooling-automation.onrender.com/}.

\subsection{User roles}
\begin{itemize}
  \item \textbf{Student.} One account does both jobs: a student can be the \emph{ride owner} of one
    ride and a \emph{rider} on another. Accounts are limited to verified IITK email addresses.
  \item \textbf{Operations administrator.} Manages hubs, pickup points, vehicle types and fare
    estimates, reviews complaints and SOS incidents, and issues warnings or suspensions. Students cannot
    open administrator routes, and administrator actions are recorded in an audit log.
\end{itemize}

\subsection{Main workflow}
\begin{enumerate}
  \item A ride owner publishes a route, departure window, vehicle type, capacity and expected fare.
  \item Riders search the eligible rides and request a seat. The owner accepts or rejects each request.
  \item The ride locks before departure, which fixes each accepted rider's share of the fare.
  \item The group coordinates in the ride chat, checks in, and records arrivals and any no-shows during
    the trip. Occupants can raise an SOS alert.
  \item When the trip completes, the system stores the final shares. Riders record their settlement,
    rate each other and can file a complaint, which an administrator reviews.
\end{enumerate}

Seats are the shared resource: several riders can compete for the last seat at the same moment, and no
seat may be allocated twice. Row locks in the request workflow and capacity triggers in the database
enforce this, and the \repolink{Deliverables/D3/Concurrency_Report/main.pdf}{concurrency report} tests it.

\subsection{AI features}
Gemini is used in three places, and in each one a person or a validated rule keeps the final decision.
\begin{itemize}
  \item \textbf{Ride recommendations} rank search results that have already been checked as eligible.
  \item \textbf{The chatbot} interprets natural-language requests and can only call allow-listed server
    operations. Changing a saved preference needs the rider's explicit confirmation.
  \item \textbf{Complaint analysis} gives an administrator advisory input. The administrator makes any
    enforcement decision.
\end{itemize}
Each feature has its own fallback if Gemini is unavailable or returns unusable output. Ride recommendations
rank the eligible rides by departure timing and fare. The chatbot interprets the request with deterministic
rules where it can, and links the rider to manual search when it cannot answer. Complaint analysis records
that no AI result is available, and the administrator reviews the complaint manually. The
\repolink{Deliverables/D3/Failure_and_Recovery_Report/main.pdf}{failure and recovery report} covers the
recommendation and chatbot cases.

\subsection{Technology}
The application is written in TypeScript with Next.js 16. A custom Node.js server hosts Socket.IO for live
ride and chat updates, and PostgreSQL 16 stores all data. Tests use Vitest and Playwright, and GitHub
Actions runs them on every pull request. It is packaged with Docker and deployed on Render.com
(Section~\ref{sec:deployment}).

\subsection{Final scope}
The requirements and design are in the
\repolink{Deliverables/D1}{Deliverable 1 documents}, and the implementation baseline, API documentation
and core workflow demonstration are in \repolink{Deliverables/D2}{Deliverable 2}. The final
\repolink{Deliverables/D3/Traceability_Matrix/main.pdf}{traceability matrix} lists the
\TraceRequirements{} Deliverable 1 requirements that a passing automated test verifies.

One decision changed after Deliverable 1: the deployment platform moved from GCP to Render.com. The
submitted requirement is unchanged, and the change is recorded in the deployment documentation
(Section~\ref{sec:deployment}).
